% title-en: Salutation to the Triple Gem
\section*{e14 | Ratanattayanāmakārapātha}

\begin{chantsubsection}
  % The Buddha
  \chantnote
    {Arahaṃ sammāsambuddho bhagavā, buddhaṃ bhagavantaṃ abhivādemi.}
    
  \chantgap
  \chantnote
    {--- bow down and chant softly ---}
  \chantnote
    {“buddho me nātho.}
  \chantnote
    {The Buddha is my refuge.”}
  \chantindent
    {(phra-phut-tha-jao pen thi-phueng khong-rao)}
  \chantgap
  \chantgap
  
  % The Dhamma
  \chantnote
    {Svākkhāto bhagavatā dhammo, dhammaṃ namassāmi.}
  \chantnote
    {--- bow down and chant softly ---}
  \chantnote
    {“dhammo me nātho.}
  \chantnote
    {The Dhamma is my refuge.”}
  \chantindent
    {(phra-dham pen thi-phueng khong-rao)}
  \chantgap
  \chantgap
  
  % The Saṅgha
  \chantnote
    {Supaṭipanno bhagavato sāvaka-saṅgho, saṅghaṃ namāmi.}
  \chantnote
    {--- bow down and chant softly ---}
  \chantnote
    {“saṅgho me nātho.}
  \chantnote
    {The Saṅgha is my refuge.”}
  \chantindent
    {(phra-song pen thi-phueng khong-rao)}
\end{chantsubsection}

% 结尾说明（可选）
\chantgap
\chantgap
{\englishstyle
This saying shows respect for the Triple Gem, firms faith in it, and keeps a close relationship with it in mind. It serves as a preliminary step that will make the attainment of Dhammakaya (the Body of Truth) more accessible.}

\chantgap
\chantgap
{\englishstyle\textit{End of the evening chanting}
}
